Термоупругость занимается вопросами равновесия тела как термодинамической системы, взаимодействие которой с окружающей средой заключается лишь в механической работе внешних сил и теплообмене. Этот раздел механики твёрдого деформируемого тела обобщает классическую теорию упругости для случаев деформаций при неравномерном нагреве \cite{Wikipedia}. Фундаментальная особенность термоупругих процессов состоит в двусторонней связи тепловых и механических явлений. Изменение температуры тела происходит не только вследствие подвода тепла от внешних источников, но и в результате самого процесса деформирования \cite{Doklad}. Эта взаимосвязь проявляется особенно ярко при быстропротекающих процессах, когда возникает так называемый эффект связанности полей деформации и температуры.

Связанность термоупругих полей обусловлена несколькими физическими механизмами. При быстрой деформации упругого тела происходит выделение или поглощение тепла, связанное с изменением объёма среды. Математически это описывается членом
\[
\gamma T_0 \frac{\partial \epsilon_{kk}}{\partial t}
\]
в уравнении теплопереноса. Неравномерный нагрев создаёт градиенты температуры, приводящие к различному тепловому расширению отдельных участков тела. Это вызывает дополнительные механические напряжения, описываемые членом
\[
- \gamma \delta_{ij} T
\]
в определяющих соотношениях. При деформировании от механических или тепловых воздействий с большой скоростью образуются и распространяются тепловые потоки внутри тела, возникают связанные волны \cite{Doklad}. Часть механической энергии волны преобразуется в тепло из-за необратимости процессов теплопроводности, что приводит к затуханию волн.

В зависимости от степени учёта взаимного влияния полей различают несвязанные и связанные задачи. Несвязанные задачи (классическая теория Дюамеля-Неймана, 1838–1841) предполагают, что температурное поле определяется независимо из уравнения теплопроводности, затем используется для расчёта температурных напряжений, и деформация не влияет на температуру. В связанных задачах учитывается обратный эффект: термоупругая деформация тела, возникающая от нестационарных механических и тепловых воздействий, сопровождается изменением его температурного поля \cite{Gemo}. Требуется одновременное решение связанной системы уравнений движения и теплопереноса.

Геологические среды характеризуются сложной структурой и неоднородностью свойств. Слоистость — чередование пород различной плотности, упругости и теплопроводности создаёт границы раздела, на которых происходит отражение и преломление волн. Анизотропия — направленная структура горных пород (например, сланцеватость, трещиноватость) приводит к зависимости скорости распространения волн от направления. Градиенты температуры, такие как геотермальный градиент и локальные тепловые аномалии, влияют на скорости волн и механические свойства среды. Пористость и насыщение — присутствие флюидов в порах усложняет термоупругое поведение за счёт дополнительных механизмов диссипации энергии.

Для тел микронных размеров амплитуда температурных волн, образующихся вследствие взаимовлияния теплового и механического полей, заметно увеличивается по сравнению с макротелами \cite{Permsc}, что имеет значение при моделировании процессов в трещиноватых и высокопористых средах.

Теория термоупругости получила существенное развитие в связи с важными проблемами, возникающими при разработке новых конструкций паровых и газовых турбин, реактивных и ракетных двигателей, высокоскоростных самолётов, ядерных реакторов \cite{Doklad}. В последние десятилетия методы термоупругости активно применяются в геофизике для моделирования сейсмических процессов, геотермальной энергетики, мониторинга подземных хранилищ и прогнозирования геомеханических рисков.